\section{Conclusiones}
La implementación del procesador RISC-V sobre FPGA permitió comprobar en la práctica los principales conceptos estudiados en arquitectura de computadoras, particularmente aquellos vinculados con segmentación de instrucciones, resolución de riesgos y control de ejecución. El desarrollo de una CPU RV32I segmentada en 5 etapas no solo implicó diseñar la ruta de datos y las señales de control, sino también enfrentar problemas reales de integración, temporización y depuración que no siempre resultan evidentes en un análisis puramente teórico.

Uno de los aportes más importantes del trabajo fue la incorporación de una \textit{Debug Unit} capaz de interactuar con la CPU mediante UART, permitiendo cargar programas, ejecutar en modo continuo o paso a paso, detener la ejecución y obtener el estado interno del sistema. Esta herramienta resultó fundamental para observar el comportamiento real del procesador, validar hipótesis durante la depuración y detectar errores que hubieran sido mucho más difíciles de aislar sin una interfaz de inspección dedicada.

Asimismo, la implementación de mecanismos de \textit{forwarding}, \textit{stall} y \textit{flush} permitió resolver correctamente los principales riesgos de datos y de control presentes en una arquitectura segmentada. Las pruebas realizadas mostraron que el procesador responde adecuadamente ante dependencias entre instrucciones, conflictos del tipo \textit{load-use} y situaciones de salto condicional o incondicional, manteniendo la correctitud de la ejecución.

A lo largo del desarrollo también surgieron dificultades concretas que enriquecieron el aprendizaje del proyecto. Entre ellas pueden mencionarse la necesidad de ajustar la latencia efectiva de las memorias para que coincidiera con la esperada por el pipeline, la depuración del comportamiento del \texttt{HALT} para lograr un vaciado real y seguro del pipeline, y la validación del protocolo de comunicación UART en presencia de transferencias extensas como los volcados de registros o memoria. Resolver estos puntos fue clave para lograr una implementación robusta y coherente.

En conjunto, los resultados obtenidos muestran que el sistema desarrollado cumple con los objetivos propuestos y constituye una implementación sólida y funcional del procesador requerido. Más allá del cumplimiento del trabajo práctico, el proyecto permitió integrar conocimientos de diseño digital, arquitectura de computadoras, verificación funcional, temporización y comunicación entre hardware y software, ofreciendo una experiencia completa y muy cercana a un desarrollo real sobre FPGA.